\section{It counts badly}
\label{sec:22.5}
Caregiving answers the identified person and fades with number. Slovic, borrowing Lifton's term, calls it \emph{psychic numbing}: feeling is greatest for a single victim and collapses as the number grows, and in later work reported compassion and hypothetical donations began to fall as soon as the count reached two, with real donations falling less clearly \autocite{slovic2007mass,vastfjall2014compassion}. The identified-victim effect is its near relation \autocite{kogut2005identified}. In a targeting pipeline the spotlight lands on the wrong side with no manipulation at all. The operator's people arrive named, with units and families. The people in the building arrive as coordinates and an estimate.

Caregiving also overrides fairness for the person it attends to. Batson and colleagues induced empathy for a child waiting for treatment, and participants who felt it were more willing to move her up the waiting list ahead of children placed before her \autocite{batson1995immorality}. For a floor, that is the vivid case defeating a term that protects unidentified people: a named hostage, a sympathetic applicant, a patrol with faces.

The scale and the weighing were set in formation for this case. The severity scale puts one family and two hundred on a ranking that doesn't depend on how either is presented. The weighing answers favoritism in part, since a vivid plight is a harm of refusing, weighed at its severity and not its vividness. Both are the reasoning Bloom recommends in place of empathy, and the construction needs them because caregiving alone would fail the way compressed review does.
